\documentclass[11pt]{article}

\usepackage[letterpaper,margin=0.85in]{geometry}
\usepackage[T1]{fontenc}
\usepackage[utf8]{inputenc}
\usepackage{array}
\usepackage{booktabs}
\usepackage{longtable}
\usepackage{tabularx}
\usepackage{xcolor}
\usepackage{hyperref}

\hypersetup{
  colorlinks=true,
  linkcolor=blue,
  urlcolor=blue
}

\setlength{\parindent}{0pt}
\setlength{\parskip}{6pt}
\renewcommand{\arraystretch}{1.15}

\newcommand{\coursename}{PH 142: Introduction to Probability and Statistics in Biology and Public Health}
\newcommand{\term}{Summer 2026}
\newcommand{\website}{https://ph142-ucb.github.io/su26/}
\newcommand{\courseemail}{jxqing@berkeley.edu}

\begin{document}

\begin{center}
{\Large \textbf{\coursename}}\\
{\large \term}\\
Summer Session D, July 6--August 14, 2026\\
4 semester credits
\end{center}

\textbf{Please note:} This syllabus is subject to change. The course website
will contain the most current schedule, links, and assignment information:
\href{\website}{\website}.

\section*{Course Format and Meeting Information}

\begin{tabularx}{\textwidth}{>{\bfseries}p{1.55in}X}
Course & PBHLTH 142, Section 001, Lecture 001\\
Class number & 13581\\
Instruction mode & Online\\
Lecture & Monday through Friday, 9:30--10:59 AM Pacific, Internet/Online. Recordings will be posted daily; live Zoom lectures will be held every other day.\\
Labs & Online and synchronous, Monday through Thursday\\
Discussion platform & Ed Discussion. We will use Ed Discussion, not Piazza, for course questions and announcements.\\
\end{tabularx}

\subsection*{Lab Sections}

\begin{tabularx}{\textwidth}{llllX}
\toprule
Section & Class \# & Days & Time & Place\\
\midrule
101A & 13583 & Mo, Tu, We, Th & 11:00--11:59 AM & Internet/Online\\
102A & 13584 & Mo, Tu, We, Th & 11:00--11:59 AM & Internet/Online\\
103A & 13585 & Mo, Tu, We, Th & 5:00--5:59 PM & Internet/Online\\
104A & 13586 & Mo, Tu, We, Th & 5:00--5:59 PM & Internet/Online\\
\bottomrule
\end{tabularx}

\section*{Course Description}

This course is an introduction to statistics and data science, primarily for
MPH and undergraduate public health majors, and others interested in public
health topics. The course can be divided into three parts.

\textbf{Part I} focuses on using R to explore datasets and summarize univariate
and bivariate distributions visually and numerically. We will also introduce
concepts around sampling and how study data are acquired. Specifically, we will
use the \texttt{dplyr} and \texttt{ggplot2} packages to manipulate and visualize
data sets in R.

\textbf{Part II} introduces classical problems in probability and the Normal,
Binomial, and Poisson distributions. We also introduce concepts around sampling
variability and the Central Limit Theorem.

\textbf{Part III} introduces statistical inference: the process of estimating
statistics from samples to make inference about populations. During all parts
of the course we will use real and simulated data sets to gain experience
conducting biostatistical analyses using R. We will follow the PPDAC model:
Problem, Plan, Data, Analysis, and Conclusion.

\section*{Prerequisites}

No specific coursework is required, and no programming background is necessary.
We expect a level of math equivalent to a US secondary school algebra course.

\section*{Course Learning Objectives}

After successfully completing Part I of the course, you will be able to:
\begin{itemize}
  \item Extract relevant statistical information from published articles in the scientific and popular press.
  \item Describe distributions of variables visually and calculate summary statistics for centrality and spread.
  \item Determine appropriate graphics for distributions and provide R code to manipulate and visualize data frames.
  \item Identify basic sampling strategies and study designs used in public health.
  \item Describe core concepts of ethics in public health.
  \item Perform basic data manipulation in R.
\end{itemize}

After successfully completing Part II of the course, you will be able to:
\begin{itemize}
  \item Compute probabilities using general probability rules.
  \item Understand statistical independence.
  \item Identify and describe Binomial and Poisson random variables.
  \item Compute probabilities using basic properties of the Normal distribution.
  \item Express epidemiologic measures as probabilities.
  \item Describe sampling variability and the Central Limit Theorem.
  \item Write R code to compute probabilities for the Normal, Binomial, and Poisson distributions.
\end{itemize}

After successfully completing Part III of the course, you will be able to:
\begin{itemize}
  \item Estimate means, proportions, and differences between means and proportions; compute confidence intervals; and perform statistical tests.
  \item Choose appropriate measures and statistical tests for associations between variables.
  \item State the assumptions and importance of assumptions for statistical tests.
  \item Describe and check assumptions for simple linear regression.
  \item Interpret confidence intervals and statistical tests for regression intercept and slope coefficients.
  \item Write R code snippets to generate confidence intervals, perform hypothesis tests, and calculate p-values.
  \item Formulate a clear statistical question, address it with an appropriate method, and accurately report findings.
\end{itemize}

\section*{Instructor Information and Communication}

\textbf{Instructor:} Jiaxin Qing (he/him/his)\\
\textbf{Instructor email:} \href{mailto:jxqing@berkeley.edu}{jxqing@berkeley.edu}\\
\textbf{Instructor availability:} Please book through email.\\
\textbf{Course email for non-content inquiries:} \href{mailto:142gsi@berkeley.edu}{142gsi@berkeley.edu}

While the instructor will interact with the whole class and oversee activities
and grading, Graduate Student Instructors (GSIs) will be your main point of
contact for assignment and course requirement questions. GSIs will lead lab
sections, hold office hours, and facilitate discussion.

All questions about course material should be asked on Ed Discussion. We will
not answer course-content questions by email. This keeps answers available to
the full class. For questions and concerns not related to course content, email
\href{mailto:142gsi@berkeley.edu}{142gsi@berkeley.edu}. Do not use the bCourses emailing
system for course communication.

\subsection*{Office Hours}

GSI office hours will be posted on the course calendar. You may attend the
office hours of any GSI regardless of your lab section. Instructor office hours
will primarily be offered by appointment through email.

\section*{Students with Disabilities}

If you are eligible for disability-related accommodations through UC Berkeley's
Disabled Students' Program, please send us your letter of accommodation through
the DSP student portal as soon as possible, so we can ensure your accommodations 
are met.

If you have any questions about your accommodations for PH142, 
please email \href{mailto:142gsi@berkeley.edu}{142gsi@berkeley.edu}. Do not 
directly email any medical documentation to the PH142 teaching team. If your 
accommodation status changes at any point during the semester, please send us 
an updated letter of accommodation through the DSP student portal.

If you believe you may be eligible for disability-related accommodations, please visit
the Disabled Students' Program website to begin an application as soon as possible. 
The application process may take several weeks.

\section*{Course Materials and Technical Requirements}

\textbf{Course website:} \href{\website}{\website}. The course website will
contain the syllabus, assignment descriptions, course resources, the most 
up-to-date schedule, and assignment information.

We will use R, RStudio, and Datahub. Use of R, RStudio, and Datahub is required
for homework assignments and lab exercises and requires an internet connection
and web browser. You will learn how to use R, RStudio, and Datahub during the
first week of classes.

\textbf{Optional textbook:} \emph{The Practice of Statistics in the Life
Sciences} by Brigitte Baldi and David S. Moore. The 4th edition is the latest
one, but previous editions are fine. The textbook is not required; course
materials take precedence where there are differences.

\section*{Learning Activities and Assignments}

All times listed are Pacific Time.

\subsection*{Lectures}

Lectures meet Monday through Friday from 9:30 to 10:59 AM Pacific. Lectures are
online. Lecture recordings will be made available every day. Live Zoom lectures
will be held every other day; on non-live days, students should use the posted
recordings and materials to keep pace with the course.

\subsection*{Lab/Discussion Sections}

Lab/discussion sections meet synchronously online Monday through Thursday.
Sections 101A and 102A meet 11:00--11:59 AM; sections 103A and 104A meet
5:00--5:59 PM. Lab attendance will be recorded, and unannnounced participation 
activities during lab sections will contribute to the participation grade.
Lab sections will include hands-on coding assignments in R, review of key concepts,
group activities, and support for group data projects.

\subsection*{Participation}

The participation grade will include unannounced lab participation activites 
and three data project GSI check-ins, in addition to other course participation
activities announced by the instructional team.

\subsection*{Homework Assignments / Problem Sets}

Homework problem sets are optional and not submitted for credit. These assignments
provide practice and preparation for exams, labs, and the data project. Solutions 
will be posted after students have had time to complete the assignment.

\subsection*{Lab Assignments}

Lab assignments are intended to practice lecture concepts in a practical
programming environment. They are intended to be completed during lab section 
with support from the GSIs. Lab exercises are submitted through Datahub to
Gradescope and are graded for correct completion. Each day's lab assignment 
will be released on that day and is due by 11:59PM in Gradescope. 

\subsection*{Midterms}

There will be three midterm exams throughout the term. The midterms will take 
place on July 16th, July 30th, and August 13th. The midterm exams are administered
via Gradescope and will be open for 24 hours. Once opened,  you have 1 hour to 
complete the exam. Exams may cover material presented in lecture, supplemental 
videos, discussion, lab sections, and R coding syntax unless otherwise noted.

Students may use one page of handwritten or printed notes for the midterm exams.
You may not use the internet to search for answers or inform your answers during exams.
You may not use AI tools during the midterm exams. While taking an exam, you 
are prohibited from discussing the exam with anyone other than the PH 142 instructional team.
Evidence of academic misconduct may result in a zero on the exam or further
disciplinary action.

\subsection*{Data Skills Demonstration Group Project}

The purpose of the group project is to use public health or biological data
that you find or have access to and use it to demonstrate statistical concepts
learned throughout the course. Data project groups are randomly assigned by the teaching
team and will be released on Tuesday, July 7th during lab section. Students are 
required to check in with their lab GSI prior to submitting each part of the data
project for participation credit.

The project will be submitted in three parts.
The Summer 2026 deadlines are:

\begin{itemize}
  \item Part I: July 17 at 11:59 PM Pacific.
  \item Part II: July 31 at 11:59 PM Pacific.
  \item Part III: August 14 at 11:59 PM Pacific.
\end{itemize}

\section*{Grading}

The final course grade percentages are:

\begin{tabularx}{\textwidth}{>{\bfseries}p{2.2in}p{1.2in}X}
\toprule
Category & Percentage & Notes\\
\midrule
Quizzes & 10\% & There are 26 quizzes total\\
Lab Assignments & 10\% & There are 11 lab assignments total\\
Midterm I & 15\% & \\
Midterm II & 15\% & \\
Midterm III & 20\% & \\
Data project & 20\% & Submitted in three parts\\
Participation & 10\% & Includes lab participation activites and data project check-ins with your GSI\\
\bottomrule
\end{tabularx}

\section*{Tentative Course Schedule}

The schedule below is tentative. The website is the authoritative schedule.

\begin{longtable}{@{}p{0.55in}p{1.05in}p{4.75in}@{}}
\toprule
Week & Dates & Topics and major deadlines\\
\midrule
\endfirsthead
\toprule
Week & Dates & Topics and major deadlines\\
\midrule
\endhead
1 & Jul 6--10 & Course introduction; PPDAC; Datahub; working with data in R; visualization; numerical summaries; relationships between variables. Labs released daily and due the same day; Weekly Assignment 1 released.\\
2 & Jul 13--17 & Regression basics; two-way tables; samples and observational studies; study design and ethics; introduction to probability. Labs released daily and due the same day; Weekly Assignment 2 released; group confirmation due Sunday.\\
3 & Jul 20--24 & Probability rules; independence; Normal, Binomial, and Poisson distributions; sampling distributions and the Central Limit Theorem. Midterm review; Weekly Assignment 3 released; Part I due Sunday.\\
4 & Jul 27--31 & Online midterm; confidence intervals; hypothesis testing; power; from z to t tests. Weekly Assignment 4 released.\\
5 & Aug 3--7 & Comparing two means; paired t tests; ANOVA; non-parametric tests; inference for proportions; inference for linear regression. Weekly Assignment 5 released; Part II due Sunday.\\
6 & Aug 10--14 & Linear regression part II; chi-square goodness of fit; chi-square independence; final review; Weekly Assignment 6 released; online final exam; Data Project Part III due Friday.\\
\bottomrule
\end{longtable}

\section*{Due Dates, Late Work, and Regrades}

Due dates for all graded assignments will be listed on the course website and
communicated in announcements on Ed Discussion.

\textbf{Late work.} Assignments submitted within 24 hours after the due date 
will be accepted with a 50\% penalty. Extensions can be made for DSP students 
and should be requested by emailing the course email. Other extension requests 
should explain the situation and must include documentation.

\textbf{Regrades.} Regrade requests must be submitted through Gradescope within
three days after grades are released. Regrades are for correcting an error in
the application of the published rubric. If you request reconsideration of one
part of an assignment, instructors may reconsider grades on the entire
assignment. Due to the short turnaround time for final grade submissions, we
generally cannot accommodate regrade requests for Midterm III.

\section*{Academic Integrity}

You are a member of an academic community at one of the world's leading
research universities. Any test, paper, code, or report submitted by you and
bearing your name is presumed to be your own original work unless otherwise
specified. You may use words or ideas from other individuals, publications,
websites, or other sources only with proper attribution.

Students may not circulate or post course materials, including handouts, exams,
syllabi, assignments, or solutions, without written permission from the
instructor. If you are unclear about expectations for completing an assignment
or taking an exam, seek clarification from the instructor or a GSI beforehand.

\section*{Strategies for Successful Learning}

Based on past experience with this course, we recommend that you:
\begin{itemize}
  \item Attend lab consistently.
  \item Ask questions in lecture, lab, office hours, and Ed Discussion.
  \item Study in a group and take turns explaining concepts to each other.
  \item Review questions and concepts from weekly assignments and exams where you lost points.
  \item Use practice exams to become familiar with question types and to practice the material.
\end{itemize}

\section*{Take Care of Yourself}

Do your best to maintain a healthy lifestyle during this intensive summer
session by eating well, exercising, getting enough sleep, and taking time to
recharge your mental health. If you start to feel overwhelmed, be kind to
yourself and reach out for support. Seeking help is a courageous thing to do
for yourself and for those who care about you.

Support resources include emotional, physical, safety, social, and basic
wellbeing resources for students. Academic resources can be found at the
Student Learning Center and English Language Resource sites. If you notice that
a classmate may need assistance, please use campus resources or reach out to
the GSIs or instructor.

\section*{Incomplete Course Grade and Course Evaluation}

Students who have substantially completed the course but, for serious
extenuating circumstances, are unable to complete remaining course activities
may request an Incomplete grade. This request must be submitted in writing to
the GSI and instructor and must include verifiable documentation.

Before the course ends, please participate in the course evaluation. Your
feedback helps improve the course and future instruction.

\end{document}
